\section{The Flat Sponge Lift}
\label{app:flat-sponge-lift}

This appendix supplies the flat sponge lift machinery
(\Cref{lem:lift-bound,lem:derived-adversary,cor:lift-application}) invoked
by the AEAD security bounds of \Cref{sec:aead-security}, the root tag bounds of
\Cref{sec:root-tag-hash-security}, and the committing-notion bound of
\Cref{cor:committing-combined}. The machinery applies uniformly to the games
named once here.

The lift has two steps in the win-event games and three in the
real-vs-ideal games. First, the \cite{BDPV08} indifferentiability theorem
replaces the real permutation used by the \TW{} construction and by direct
primitive queries with the simulator over a flat random oracle, at cost
$\Lift_c(\Ntot)$, where $\Ntot = N + \Nprim$ bounds the operational primitive
trace: construction-internal duplex calls induced by construction queries and
deterministic challenger computations, together with direct primitive queries.
The simulator-side construction trace is then identified with the flat-RO
construction interface by the operational trace lemma below.
Second, the simulator-world adversary is converted into a flat-RO
adversary whose direct random oracle query count is at most $\Nprim$. Thus the
random oracle bounds are instantiated at $\qRO \leq \Nprim$, while the public
permutation lift pays $\Lift_c(N+\Nprim)$. Third, in the real-vs-ideal games,
whose ideal world exposes the real permutation rather than the simulator
(\Cref{sec:games-pp}), the ideal-side simulator introduced by the first step
is exchanged back for $(\pi,\pi^{-1})$ at cost $\Lift_c(\Nprim)$
(\Cref{lem:sim-perm-proximity}); the exchange is possible because the ideal
construction oracles are independent of the primitive.

\begin{definition}[Single-Stage \TW{} Games]
\label{def:single-stage-games}
The \emph{single-stage \TW{} games} are the games of
\Cref{sec:games-pp,sec:mu-setting} in which one adversary interacts with the
challenger in a single stage, so that a \cite{BDPV08} replacement applies to
an entire world's public-permutation transcript at once, including
deterministic challenger checks after the adversary's final output. They are
all-in-one AE, mu-ind, \textsf{INT-RUP}, CMTDs-4, the $\HTW$ collision and
preimage games (standard and everywhere), and CMTs-$\ell$ for
$\ell \in \{1,3,4\}$. The win-event games among these have one world; the
real-vs-ideal games---all-in-one AE and mu-ind---have two, and the lift below
treats their real world by the full replacement and their ideal world by the
simulator--permutation exchange of \Cref{lem:sim-perm-proximity}.
\end{definition}

The derived adversary forwards the following game data to the corresponding
flat-RO experiment.
\begin{center}
\small
\begin{tabular}{@{}l p{0.62\linewidth}@{}}
  \toprule
  Game family & Forwarded data \\
  \midrule
  AE / mu-ind / \textsf{INT-RUP}
    & construction queries, including $\mathsf{New}$ where present and replay state \\
  $\HTW$ hash games & final structured hash game output \\
  CMTDs / CMTs & final openings and deterministic challenger checks \\
  \bottomrule
\end{tabular}
\end{center}

\subsection{Simulator Convention}
\label{app:simulator}

Throughout this appendix, $(\Sim^{\ROflat}, \Sim^{-1,\ROflat})$
denotes the public primitive simulator supplied by \cite[Theorem~2]{BDPV08}
--- the indifferentiability theorem for the simple
padded sponge over a random permutation --- after applying the
\cite[Theorem~3]{BDPV11} bridge of \Cref{lem:bridge-decode} to \TW{}'s native
$\padtenone$ transcript inputs. The superscript records that both
simulator directions share the same internal lazy $\ROflat$ table.

The simulator is a proof device: each game of
\Cref{sec:games-pp,sec:mu-setting} states its advantage over the real
permutation, and the reduction to the corresponding random oracle model
adversary passes through intermediate simulator worlds. In those worlds,
the construction interface is the one specified by the game
($(\Rand,\Replay)$, instantiated per user in mu-ind, on the ideal side)
and the public primitive interface is
$(\Sim^{\ROflat}, \Sim^{-1,\ROflat})$. The $\ROflat$ table in this
notation is internal to that system and is not an extra interface
exposed to the application adversary.

\begin{lemma}[At Most One Simulator RO Call]
\label{lem:sim-ro-call-filter}
For the simulator $(\Sim^{\ROflat},\Sim^{-1,\ROflat})$, every
primitive-interface query induces at most one simulator-side $\ROflat$ call.
More precisely, every inverse primitive-interface query and every forward
primitive-interface query violating one of the following conditions induces no
simulator $\ROflat$ call. The symbols $s$, $p$, $\mathcal{R}$, $\mathcal{O}$,
and $\mathcal{C}$ are the local internal notation of the BDPV simple-padded
sponge simulator, not \TW{} resource notation:
\begin{description}
  \item[(C1) New edge.] The simulator's internal graph has no
    outgoing edge from the input node $s$.
  \item[(C2) Unsaturated.] $\mathcal{R} \cup \mathcal{O} \neq
    \mathcal{C}$.
  \item[(C3) Rooted input.] The input node is rooted.
  \item[(C4) Paddable path.] The constructed path $p$ decomposes
    uniquely as $p = p' \concat 0^{rj}$ with $p'$ not ending in
    $0^r$, and the prefix $p'$ is a valid padded sponge input:
    equivalently, $p' = x \concat \padten[r](|x|)$ for an unpadded
    BDPV08 random oracle input $x$.
\end{description}
When all four conditions hold, the forward primitive-interface query
induces one simulator $\ROflat$ call. Consequently, each
primitive-interface query induces at most one simulator-side
$\ROflat$ call.
\end{lemma}

\begin{proof}
This is the branch structure of the \cite{BDPV08} simple-padded
permutation simulator. The simulator invokes the random oracle only on
the forward-query branch where the input
node is rooted, unsaturated, paddable to an unpadded random oracle
input $x$, and not already extended by an outgoing edge in the
simulator graph. Forward queries violating any
of these conditions and all inverse queries return outputs sampled or
read from the simulator's internal graph without invoking the
random oracle. Under the \cite[Theorem~3]{BDPV11} bridge fixed in
\Cref{app:simulator}, that simulator random oracle invocation is the
simulator-side $\ROflat$ call used throughout this appendix.
\end{proof}

\begin{lemma}[Operational Duplex Trace]
\label{lem:operational-duplex-trace}
In any execution of a single-stage \TW{} game
(\Cref{def:single-stage-games}), the construction computations counted
by $N$ admit an operational primitive trace of cost at most $N$:
each root or leaf duplex call contributes one forward primitive query,
and the requested output bits are the corresponding projection of the
state returned by that query. When this trace is run against the
\cite{BDPV08} simulator $(\Sim^{\ROflat},\Sim^{-1,\ROflat})$ in an
execution whose combined primitive trace has cost less than $2^c$, the
construction outputs are the flat-RO values
$\RF(X)=\ROflat(\flatmap(X))$ at the same well-formed \TW{} transcript
prefixes $X$.
\end{lemma}

\begin{proof}
A \TW{} duplex call absorbs a single suffixed input block, padded by
$\padtenone$, applies the permutation once, and returns at most $r$
rate bits from the resulting state. Thus the public-permutation
implementation of any root or leaf construction computation is already
a primitive trace with one forward query per duplex call. The definition
of $N$ in \Cref{sec:resources} counts exactly these root initialization,
associated data, message, leaf initialization, leaf message, and
aggregation duplex calls, including the deterministic challenger checks
of the hash and committing games.

\Cref{lem:duplex-flat} is the \cite[Lemma~3]{BDPV11} duplex-to-sponge
identity specialized to \TW{}: the output of each duplex call is the
output of the sponge after absorbing the accumulated flattened
duplex-call transcript, and no additional permutation call is made after
the current absorption block when the requested output length is at most
$r$. The bridge of \Cref{lem:bridge-decode}, fixed in
\Cref{app:simulator}, maps these $\padtenone$ transcript prefixes to the
\cite{BDPV08} padded-sponge input domain without changing the output bits
at $r=r_{\max}$.

When the same operational trace is run against
$(\Sim^{\ROflat},\Sim^{-1,\ROflat})$, every construction-internal forward
query lies on a rooted, paddable path determined by a well-formed \TW{}
prefix. In an execution whose combined primitive trace has cost less
than $2^c$, the simulator cannot have saturated: each primitive query
adds at most one capacity value to the simulator's internal
$\mathcal{R} \cup \mathcal{O}$ set, so fewer than $2^c$ queries cannot
fill the $2^c$-element capacity space. The simulator's
sponge-consistency invariant \cite[Lemma~2]{BDPV08} therefore makes the
rate part returned at that node agree with the corresponding lazy value
$\ROflat(\flatmap(X))$. Thus the simulator-side operational construction
has the same construction-output law as the flat-RO construction
interface used in \Cref{sec:ro-model}.
\end{proof}

\begin{lemma}[Flat Sponge Lift Bound]
\label{lem:lift-bound}
Let $\mathsf{G}$ be a single-stage \TW{} game
(\Cref{def:single-stage-games}). Let
$\mathcal{A}$ be a
public permutation $\mathsf{G}$ adversary with construction interface
cost at most $N$ and at most $\Nprim$ direct primitive-interface
queries. In the \cite{BDPV08} single-stage replacement, implement the
\TW{} construction computations by the operational duplex trace of
\Cref{lem:operational-duplex-trace} and keep the adversary's direct
primitive-interface queries as primitive queries. This primitive
transcript has \cite{BDPV08} query cost at most
\[
  \Ntot = N + \Nprim .
\]
Consequently replacing the real public permutation and \TW{} construction
by the simulator-world primitive interface
$(\Sim^{\ROflat},\Sim^{-1,\ROflat})$ and the corresponding flat-RO
construction interface costs at most $\Lift_c(\Ntot)$.
\end{lemma}

\begin{proof}
The \cite{BDPV08} replacement is applied before the derived-adversary
construction. View the entire single-stage game, including its
construction interface and deterministic challenger checks, as a
primitive-only distinguisher that runs the \TW{} algorithms internally.
By \Cref{lem:operational-duplex-trace}, the construction side contributes
at most $N$ forward primitive queries: one per duplex call. This includes
plaintext-computing computations, verification computations whether they
accept or reject, the per-user sum of construction calls in mu-ind, the
deterministic challenger decryption or encryption checks charged in the
committing games, and the hash games' challenger encryptions---the
collision game's $s$ submitted-input encryptions, the standard preimage
game's source encryption and post-output check, and the
everywhere-preimage game's post-output check.

Each direct adversary query to $(\pi,\pi^{-1})$ contributes one
primitive-interface query and is counted in $\Nprim$. These two counts
are disjoint in the public permutation experiment: $N$ counts
construction-internal use of the primitive, while $\Nprim$ counts only
the adversary's public primitive interface queries.

Construction work is modeled by the operational duplex trace, rather
than by a collection of standalone \cite{BDPV08} construction queries at
every flattened duplex prefix. The standalone-query model would assign a
nested-prefix transcript the sum of the individual sponge-evaluation
costs, while the operational trace is the same stateful computation the
construction actually performs, with the output of each call read from
the state just returned by that call. Hence the execution-uniform caps
bound the primitive-only query cost by $N+\Nprim=\Ntot$. If
$\Ntot \geq 2^c$ then $\Lift_c(\Ntot) \geq 1$ and the claimed bound holds
trivially, so assume $\Ntot < 2^c$, satisfying the hypothesis of
\cite[Theorem~2]{BDPV08}. Applying that theorem to this primitive-only
distinguisher gives the real-to-simulator replacement, and
\Cref{lem:operational-duplex-trace} identifies the simulator-side
construction outputs with the flat-RO construction interface in this
nontrivial case. The theorem gives distinguishing or
success-probability cost at most the random-permutation
differentiating bound $f_P(\Ntot)$ of \cite[Lemma~5]{BDPV08}. The
\cite[Lemmas~4 and~5]{BDPV08} product bounds give the
random-transformation bound $f_T(\Ntot) = 1 - \prod_{i=1}^{\Ntot}(1 -
i/2^c)$, and the factor of $f_P(\Ntot)$ corresponding to the $f_T$ factor
at index $i$ is that factor divided by a quantity
$1 - (i-1)/2^{r+c} \in (0, 1]$, hence at least as large, so
$f_P(\Ntot) \leq f_T(\Ntot)$. With $\Ntot < 2^c$ every factor lies in
$[0, 1]$ and
\[
  f_P(\Ntot) \leq f_T(\Ntot)
  = 1 - \prod_{i=1}^{\Ntot}\Bigl(1 - \frac{i}{2^c}\Bigr)
  \leq \sum_{i=1}^{\Ntot} \frac{i}{2^c}
  = \frac{\Ntot(\Ntot + 1)}{2^{c+1}}
  = \Lift_c(\Ntot),
\]
the second inequality by the Weierstrass product inequality
$\prod_i (1 - x_i) \geq 1 - \sum_i x_i$ for $x_i \in [0, 1]$.
\end{proof}

\begin{lemma}[Simulator--Permutation Proximity]
\label{lem:sim-perm-proximity}
Let $\mathcal{D}$ be a distinguisher with access only to a two-directional
primitive interface, making at most $q$ queries. Then
\[
  \bigl|
    \Pr[\mathcal{D}^{\Sim^{\ROflat},\, \Sim^{-1,\ROflat}} \Rightarrow 1]
    - \Pr[\mathcal{D}^{\pi,\, \pi^{-1}} \Rightarrow 1]
  \bigr|
  \;\leq\; \Lift_c(q),
\]
where $\pi \sample \Perm(\bits^{1600})$ and $\ROflat$ is the lazily
sampled random oracle internal to the simulator.
\end{lemma}

\begin{proof}
Consider the indifferentiability distinguisher $\mathcal{D}'$ that runs
$\mathcal{D}$, forwards its primitive-interface queries, and never queries
the construction interface. The \cite{BDPV08} query cost of $\mathcal{D}'$
is at most $q$: each forwarded query costs one, and no construction queries
arise. Because $\mathcal{D}'$ ignores the construction interface, its two
views are exactly $\mathcal{D}^{\pi,\pi^{-1}}$ in the real pair and
$\mathcal{D}^{\Sim^{\ROflat},\Sim^{-1,\ROflat}}$ in the ideal pair of
\cite[Theorem~2]{BDPV08}, applied under the \cite[Theorem~3]{BDPV11}
bridge fixed in \Cref{app:simulator}. If $q \geq 2^c$ then
$\Lift_c(q) \geq 1$ and the claim is trivial; otherwise that theorem
bounds the displayed difference by the random-permutation differentiating
bound $f_P(q)$ of \cite[Lemma~5]{BDPV08}, and the derivation in the proof
of \Cref{lem:lift-bound} gives $f_P(q) \leq \Lift_c(q)$.
\end{proof}

\begin{lemma}[Derived-Adversary Execution Contract]
\label{lem:derived-adversary}
Let $\mathsf{G}$ be a single-stage \TW{} game
(\Cref{def:single-stage-games}). Let
$\mathcal{A}$ be a
simulator-world $\mathsf{G}$ adversary whose public primitive
interface is answered by
$(\Sim^{\ROflat}, \Sim^{-1,\ROflat})$, with at most $\Nprim$ direct
primitive-interface queries. Define
$\mathcal{B}_{\mathrm{derived}}(\mathcal{A})$ to run $\mathcal{A}$
internally, answer each primitive-interface query by running the
same simulator against direct finite-output queries to its own
$\ROflat$ oracle, forward construction interface queries unchanged in games
with construction oracles, and forward the adversary's final structured output
unchanged in the output-checked $\HTW$ hash and committing games. Then the
simulator-world execution of $\mathcal{A}$ and the random oracle execution of
$\mathcal{B}_{\mathrm{derived}}(\mathcal{A})$ have identical joint
distributions, and
\[
  \qRO(\mathcal{B}_{\mathrm{derived}}) \leq \Nprim.
\]
In this coupling, simulator-induced $\ROflat$ calls and
construction-internal $\RF$ lookups use the same lazy table. Therefore,
if a simulator-induced query is the bridged input $\flatmap(X)$ of a
well-formed \TW{} transcript prefix $X$ and a forwarded construction
computation later reaches the same prefix $X$, both interfaces read the
same $\ROflat(\flatmap(X))$ value, projected to the requested output
length.
Moreover, $\mathcal{B}_{\mathrm{derived}}$ inherits the construction-side
caps of $\mathcal{A}$:
\begin{align*}
  \qv(\mathcal{B}_{\mathrm{derived}}) &= \qv(\mathcal{A})
    \quad\text{(AE, mu-ind, \textsf{INT-RUP})}, \\
  \nleaf(\mathcal{B}_{\mathrm{derived}}) &\leq \nleaf(\mathcal{A})
    \quad\text{(all games above)}.
\end{align*}
In the mu-ind game these resource inheritances hold per user, and the
$\mathsf{New}$-query transcript is unchanged.
\end{lemma}

\begin{proof}
Couple the two executions by using the same lazy $\ROflat$ table for
the simulator and for all construction-internal $\RF(\cdot)$ lookups.
In the simulator-world execution, $\mathcal{A}$ sees construction
answers from the selected game interface and public primitive answers
from $(\Sim^{\ROflat}, \Sim^{-1,\ROflat})$. In the derived
random oracle execution, $\mathcal{B}_{\mathrm{derived}}$ runs the
same $\mathcal{A}$, forwards construction interface queries or the final
structured output required by the selected game unchanged, and answers
primitive-interface queries by invoking the same simulator code with the same
lazy $\ROflat$ table. The construction interface, hash challenger, or
committing challenger therefore receives the same forwarded queries or final
output, the simulator receives the same primitive queries in the same order, and
the lazy table supplies the same values.
The full views of $\mathcal{A}$ in the simulator-world execution and of
$\mathcal{B}_{\mathrm{derived}}$ in the random oracle execution are identically
distributed.

The simulator's random oracle calls may be on arbitrary
\cite{BDPV08} $H$-input strings, not only bridged \TW{} transcript
inputs. These calls are the direct $\ROflat$ queries of
$\mathcal{B}_{\mathrm{derived}}$ and are counted in
$\qRO(\mathcal{B}_{\mathrm{derived}})$. Construction-internal
$\RF(\cdot)$ lookups made by $\Enc_K^\RO$, $\Dec_K^\RO$,
$\PDec_K^\RO$, $\Ver_K^\RO$, or a committing challenger are forwarded to the
same $\ROflat$ table but are not counted in $\qRO$, exactly as in
\Cref{sec:resources}.

By \Cref{lem:sim-ro-call-filter}, each primitive-interface query
induces at most one direct $\ROflat$ query by
$\mathcal{B}_{\mathrm{derived}}$, and inverse queries or
non-RO-touching forward queries induce none. Since $\mathcal{A}$ makes
at most $\Nprim$ direct primitive-interface queries, this gives
$\qRO(\mathcal{B}_{\mathrm{derived}}) \leq \Nprim$.

\Cref{lem:bridge-decode}'s bridge places \TW{} construction
restrictions and simulator-induced $\ROflat$ calls in a single lazy
table. If the simulator invokes $\ROflat$ on $\flatmap(X)$ for a
well-formed \TW{} transcript prefix $X$, and a forwarded construction
computation later reaches $X$, the construction lookup
$\RF(X)=\ROflat(\flatmap(X))$ reads the same table entry and returns the
requested projection. Since \TW{} uses $r = r_{\max}$, the
\cite[Theorem~3]{BDPV11} output post-processing does not alter \TW{}
output bits. Thus simulator-induced direct queries and
construction-internal lookups at the same bridged input are
transcript-consistent. In the AE-style games, a simulator call whose input is
accepted by $\KeyedTWOut$ for the hidden key is still a direct
$\ROflat$ query and is covered by the same
$\badkey$ accounting used by the direct random oracle mu-ind proof. The
derivation step itself does not
prove a key-guessing statement; it only supplies the derived
adversary and the query bound at which the random oracle bounds of
\Cref{sec:aead-security,sec:root-tag-hash-security} are instantiated.

Finally, the execution-uniform caps of \Cref{sec:resources} hold
along every oracle-answer trace of $\mathcal{A}$, including traces
generated by the simulator on the ideal public primitive interface.
The derived-adversary construction does not add construction interface
encryption, plaintext-computing, or verification queries and forwards
$\mathcal{A}$'s construction interface queries unchanged along each trace.
$\mathcal{B}_{\mathrm{derived}}$ therefore inherits the relevant
construction resources: $\qv(\mathcal{B}_{\mathrm{derived}}) =
\qv(\mathcal{A})$ in the AE, mu-ind, and \textsf{INT-RUP} games, and
$\nleaf(\mathcal{B}_{\mathrm{derived}}) \leq \nleaf(\mathcal{A})$ for
the distinct construction-generated final-leaf-transcript cap. In
mu-ind, construction interface forwarding preserves the
$\mathsf{New}$ calls and these caps user by user.
\end{proof}

\subsection{Lift Application}
\label{app:lift-corollaries}

The lift bound of \Cref{lem:lift-bound}, the simulator--permutation
proximity bound of \Cref{lem:sim-perm-proximity}, and the
derived-adversary construction of \Cref{lem:derived-adversary} combine
into a single template, stated once and then instantiated by the bounds of
\Cref{sec:aead-security,sec:root-tag-hash-security} and by the
committing-notion bounds of \Cref{cor:committing-combined}. Every
public-permutation bound is obtained by adding $\Lift_c(\Ntot)$---plus
$\Lift_c(\Nprim)$ in the real-vs-ideal games---and
substituting $\qRO \leq \Nprim$ in the corresponding random oracle theorem.

\begin{corollary}[Lift Application]
\label{cor:lift-application}
Let $\mathsf{G}$ be a single-stage \TW{} game
(\Cref{def:single-stage-games}). For every
public permutation $\mathsf{G}$-adversary $\mathcal{A}$ against \TW{}
with the resources of \Cref{sec:resources}, the derived adversary
$\mathcal{B} = \mathcal{B}_{\mathrm{derived}}(\mathcal{A})$ of
\Cref{lem:derived-adversary} is a random oracle $\mathsf{G}$-adversary
with $\qRO(\mathcal{B}) \leq \Nprim$ and the construction-side caps of
\Cref{lem:derived-adversary}---in particular $\qv(\mathcal{B}) \leq \Qv$
and $\nleaf(\mathcal{B}) \leq \Nleaf$, where these apply---such that,
when $\mathsf{G}$ is a win-event game,
\[
  \Adv^{\mathsf{G}}_\TW(\mathcal{A})
  \;\leq\;
  \Lift_c(\Ntot) + \Adv^{\mathsf{G}}_{\RO}(\mathcal{B}),
\]
and, when $\mathsf{G}$ is a real-vs-ideal game (all-in-one AE or mu-ind),
\[
  \Adv^{\mathsf{G}}_\TW(\mathcal{A})
  \;\leq\;
  \Lift_c(\Ntot) + \Adv^{\mathsf{G}}_{\RO}(\mathcal{B}) + \Lift_c(\Nprim).
\]
If moreover $\Adv^{\mathsf{G}}_{\RO}(\mathcal{B}) \leq
\varepsilon_{\mathsf{G}}(\qRO(\mathcal{B}), \nleaf(\mathcal{B}))$ for some
$\varepsilon_{\mathsf{G}}$ nondecreasing in each argument, then
$\varepsilon_{\mathsf{G}}(\Nprim, \Nleaf)$ may replace
$\Adv^{\mathsf{G}}_{\RO}(\mathcal{B})$ in the corresponding bound.
\end{corollary}

\begin{proof}
By the flat sponge lift bound (\Cref{lem:lift-bound}), replacing the real
public permutation and \TW{} construction by
$(\Sim^{\ROflat},\Sim^{-1,\ROflat})$ and the flat-RO construction
interface costs at most $\Lift_c(\Ntot)$, where $\Ntot = N + \Nprim$
counts both the construction calls of $N$---including any construction
oracle exposed during $\mathcal{A}$'s run and, in the win-event games,
the challenger's construction computations, pre- and post-output (in the
preimage game, the source encryption computing $T^*$ as well as the
post-output check)---and the
$\Nprim$ direct primitive queries. The derived-adversary contract
(\Cref{lem:derived-adversary}) expresses the resulting simulator-world
execution as the identically distributed random oracle execution of
$\mathcal{B} = \mathcal{B}_{\mathrm{derived}}(\mathcal{A})$, with
$\qRO(\mathcal{B}) \leq \Nprim$ and the stated construction-side
inheritances. For a win-event game this gives the first bound.

For a real-vs-ideal game, chain four worlds:
$\mathsf{W}_1$, the real construction interface with $(\pi,\pi^{-1})$;
$\mathsf{W}_2$, the flat-RO real construction interface with the
simulator; $\mathsf{W}_3$, the ideal construction interface with the
simulator; and $\mathsf{W}_4$, the ideal construction interface with
$(\pi,\pi^{-1})$. The game of \Cref{sec:games-pp,sec:mu-setting}
compares $\mathsf{W}_1$ with $\mathsf{W}_4$. The step from
$\mathsf{W}_1$ to $\mathsf{W}_2$ costs at most $\Lift_c(\Ntot)$ as
above, and under the coupling of \Cref{lem:derived-adversary} the pair
$(\mathsf{W}_2, \mathsf{W}_3)$ is precisely the real/ideal pair of the
random oracle game for $\mathcal{B}$, so
$|\Pr[\mathcal{A}^{\mathsf{W}_2} \Rightarrow 1] -
\Pr[\mathcal{A}^{\mathsf{W}_3} \Rightarrow 1]| \leq
\Adv^{\mathsf{G}}_{\RO}(\mathcal{B})$. For the step from $\mathsf{W}_3$
to $\mathsf{W}_4$, the ideal construction oracles---$(\Rand,\Replay)$,
instantiated per user in mu-ind, together with the $\mathsf{New}$
interface and the validity checks---are independent of the primitive, so
a distinguisher can answer them from its own coins, run $\mathcal{A}$,
and forward the at most $\Nprim$ primitive-interface queries;
\Cref{lem:sim-perm-proximity} bounds this step by $\Lift_c(\Nprim)$.
The triangle inequality combines the three estimates into the second
bound.

The final claim follows by the
monotonicity of $\varepsilon_{\mathsf{G}}$. When $\mathsf{G}$ is
parameterized by a challenge fixed before $\mathcal{A}$ runs---the
everywhere-preimage target tag, or the standard preimage source $X^*$,
whose tag $T^* = \HTW(X^*)$ each coupled world computes---the coupling
holds for each fixed challenge and the challenge-independent bound holds
for the maximum or average over challenges defining
$\Adv^{\mathsf{G}}_\TW(\mathcal{A})$.
\end{proof}
